\subsection{A Gaussian construction for regularized normalization}
\label{sec:gaussianrms}

A positive stabilizer permits a factory with output radius $4\sqrt n$
and source bound $1+nR^2/\epsilon$. The construction uses the Gaussian
sign identity, an error-function mixture of majority rules, and a sine
series. Its denominator parameter is the stabilizer itself.

Write
\[
 \operatorname{erf}(x)=\frac2{\sqrt\pi}\int_0^x e^{-t^2}\,dt,
 \qquad
 A_k=\frac{(2k+1)\binom{2k}{k}}{4^k}.
\]
For a sign of mean $z$, let $M_k(z)$ be the mean of the majority of
$2k+1$ independent copies. The derivative identity
$M'_k(z)=A_k(1-z^2)^k$ from the proof of
Lemma~\ref{lem:comptanharity} and $M_k(0)=0$ imply
\begin{equation}
 \operatorname{erf}(cz)=\sum_{k\ge0}
 \frac{2c}{\sqrt\pi}e^{-c^2}\frac{c^{2k}}{k!A_k}M_k(z),
 \qquad c\ge0,\quad |z|\le1.
 \label{eq:erfmajority}
\end{equation}
All coefficients are nonnegative, and their sum is $\operatorname{erf}(c)$.

\begin{lemma}[An error-function factory]
\label{lem:erffactory}
For every known computable $c\ge0$, there is an exact factory with
output mean $\operatorname{erf}(cz)$. Its expected number of input
signs, using complete majorities, is exactly
\[
 q(c)=2c^2\operatorname{erf}(c)+\frac{2c}{\sqrt\pi}e^{-c^2}
 \le1+2c^2.
\]
One implementation draws $G\sim N(0,1)$. If $G^2>2c^2$, it returns
a fair sign. Otherwise it draws
$K\sim\operatorname{Pois}(c^2-G^2/2)$ and returns the majority of
$2K+1$ fresh input signs.
\end{lemma}
\paragraph{Proof idea.}
The derivative of an odd-majority mean is a power of $1-z^2$.
Poisson coefficients turn their sum into the Gaussian density, and a
truncated Gaussian draw realizes those coefficients as probabilities.
The same integral computes the expected majority size.

\begin{proof}
Differentiating the right side of \eqref{eq:erfmajority}, and using
nonnegative convergence on $[0,1]$, gives
\[
 \frac{2c}{\sqrt\pi}e^{-c^2}
 \sum_{k\ge0}\frac{[c^2(1-z^2)]^k}{k!}
 =\frac{2c}{\sqrt\pi}e^{-c^2z^2}.
\]
Both sides vanish at zero; oddness handles negative $z$. At $z=1$
every majority has mean one, proving the coefficient sum.

The variable $Y=|G|/\sqrt2$ has density
$(2/\sqrt\pi)e^{-y^2}$ for $y\ge0$. For each $k$, the probability
that the algorithm enters the majority branch and selects $k$ is
\[
 \frac2{\sqrt\pi}\int_0^c e^{-y^2}
       e^{-(c^2-y^2)}\frac{(c^2-y^2)^k}{k!}\,dy
 =\frac{2c}{\sqrt\pi}e^{-c^2}\frac{c^{2k}}{k!}
       \int_0^1(1-u^2)^k\,du.
\]
Integration by parts, starting from $k=0$, gives
$\int_0^1(1-u^2)^k\,du=1/A_k$. The algorithm therefore uses the
mixture in \eqref{eq:erfmajority} exactly. It never forms a square
of an unknown mean.

The number of input signs has conditional expectation
$1+2(c^2-G^2/2)$ on the majority branch. Consequently its mean is
\[
 \frac2{\sqrt\pi}\int_0^c
   e^{-y^2}(1+2c^2-2y^2)\,dy
 =2c^2\operatorname{erf}(c)+\frac{2c}{\sqrt\pi}e^{-c^2},
\]
where the last equality uses
$(y e^{-y^2})'=e^{-y^2}-2y^2e^{-y^2}$.
The bound $1+2c^2$ also follows directly from the conditional
expectation. A sequential majority may stop earlier.
\end{proof}

\begin{lemma}[A sine factory]
\label{lem:sinequarter}
Given independent signs of unknown mean $z\in[-1,1]$, an exact
factory of mean $\frac14\sin(\pi z/2)$ uses
\[
 C_{\sin}=\frac\pi8\cosh(\pi/2)<1
\]
input signs in expectation.
\end{lemma}
\paragraph{Proof idea.}
The sine series has a summable total absolute coefficient mass
below one at the chosen scale. Products of fresh signs realize its monomials;
a fair-sign remainder supplies the unused probability mass.

\begin{proof}
Put $c_0=\pi/2$. For each $k\ge0$, choose degree $2k+1$ with
probability $c_0^{2k+1}/[4(2k+1)!]$. Return $(-1)^k$ times the
product of that many independent input signs. On the remaining
probability return a fair sign. The total specified mass is
$\sinh(c_0)/4<1$, and the Taylor series gives the stated mean.
The weighted degree sum is $c_0\cosh(c_0)/4=C_{\sin}$.

For an explicit rational check on the strict inequalities, use
$\pi<22/7$ and
\[
 \cosh(11/7)\le
 1+\frac{(11/7)^2}{2!}+\frac{(11/7)^4}{4!}
 +\frac{(11/7)^6}{6!}
 +\frac{(11/7)^8}{8![1-(11/7)^2/90]}<\frac{101}{40}.
\]
Thus $C_{\sin}<2222/2240<1$, and also
$\sinh(c_0)/4<\cosh(c_0)/4<101/160<1$.
The series can be sampled by an odd Poisson variable with parameter
$c_0$, preceded by the known gate $\sinh(c_0)/4$. All parameters
are computable constants, and the degree has an exponential tail.
\end{proof}

\begin{theorem}[An intrinsic-radius RMS factory]
\label{thm:gaussianrms}
Let $n\ge1$, $R>0$, and $\epsilon>0$. Suppose fresh independent
signs of means $a_j/R$ are available, where $|a_j|\le R$.
For every coordinate $i$, an exact factory returns a sign of mean
\[
 \frac{a_i}{4\sqrt n\sqrt{\epsilon+n^{-1}\sum_j a_j^2}}.
\]
Its expected number of source-sign requests is at most
\begin{equation}
 C_{\sin}\left(1+\frac{R^2}{\epsilon}
      \left[1+\frac{2(n-1)}\pi\right]\right)
 \le1+\frac{nR^2}{\epsilon}.
 \label{eq:gaussianrmscost}
\end{equation}
All Gaussian and Poisson variables can be represented by successively
refined rational intervals. For finite rational $R,\epsilon$, and
$B\ge16$ containing their encoding lengths and $\log(n+2)$, the
auxiliary expected bit work is
\[
 O\!\left(n(1+R^2/\epsilon)B^4\right).
\]
This counts Gaussian draws, random bits, arithmetic, and precision
refinement; it uses no exact-real oracle.
\end{theorem}
\paragraph{Proof idea.}
A Gaussian projection converts the normalized coordinate into the
correlation of two Gaussian signs. That correlation is an arcsine, which
the sine factory inverts. Fresh Gaussian vectors for successive calls let
the source costs multiply by conditional expectation.

\begin{proof}
Let $G_1,\ldots,G_n$ be independent standard Gaussian variables.
Conditional on this vector, put
\[
 A=R\sum_j|G_j|,\qquad
 c=\frac{A}{\sqrt{2n\epsilon}},\qquad
 z=\frac{\sum_jG_ja_j}{A}.
\]
The event $A=0$ has probability zero. A source sign of mean $z$ is
obtained by selecting $j$ proportionally to $|G_j|$, requesting a
fresh sign of mean $a_j/R$, and multiplying it by $\operatorname{sgn}G_j$.
Use Lemma~\ref{lem:erffactory} and multiply its output by
$\operatorname{sgn}G_i$. This produces a sign whose conditional mean is
\[
 \operatorname{sgn}G_i\,
 \operatorname{erf}\left(\frac{\sum_jG_ja_j}{\sqrt{2n\epsilon}}\right).
\]

Introduce an independent $G_0\sim N(0,1)$ only for the analysis.
Integrating it out gives the error function above as the conditional
mean of
$\operatorname{sgn}(\sum_jG_ja_j+\sqrt{n\epsilon}G_0)$.
The two centered Gaussian variables
\[
 G_i,\qquad H=\sum_jG_ja_j+\sqrt{n\epsilon}G_0
\]
have correlation
$u=a_i/\sqrt{\sum_j a_j^2+n\epsilon}$.
The Gaussian sign identity gives
\[
 \mathbb E[\operatorname{sgn}G_i\operatorname{sgn}H]
 =\frac2\pi\arcsin u.
\]
For completeness, write a pair with correlation $u=\cos\theta$ as
$(X,\cos\theta X+\sin\theta Y)$ with independent standard Gaussians
$X,Y$. Their angle is uniform on the circle; the two signs disagree
on sectors of total angle $2\theta$. Their product therefore has
mean $1-2\theta/\pi=(2/\pi)\arcsin u$.

Request independent copies of this sign as inputs to
Lemma~\ref{lem:sinequarter}. The returned mean is
\[
 \frac14\sin(\arcsin u)=\frac u4,
\]
which is the desired output mean.
Each request draws its own Gaussian vector and all subsequent
randomness afresh.

By Lemma~\ref{lem:erffactory}, a single intermediate sign uses at
most $1+2c^2$ source requests. Since
$\mathbb E G_j^2=1$ and $\mathbb E|G_j|=\sqrt{2/\pi}$,
\[
 2\mathbb E c^2
 =\frac{R^2}{n\epsilon}\mathbb E\left(\sum_j|G_j|\right)^2
 =\frac{R^2}{\epsilon}\left[1+\frac{2(n-1)}\pi\right].
\]
Multiplication by the mean input count of the independent sine
factory proves \eqref{eq:gaussianrmscost}. Its degree is drawn before
any of its input signs, so this multiplication needs only Tonelli's
theorem and conditional expectation. The finite-bit implementation
is proved in the accompanying precision lemma.
\end{proof}

The vector result provides a second option when $\epsilon>0$.
The floor factory of Theorem~\ref{thm:rmsfloor} has radius
$2R/\sqrt\lambda$ and cost $1+228\max\{1,R^2/\lambda\}^2$.
The Gaussian factory has radius $4\sqrt n$ and cost at most
$1+nR^2/\epsilon$. A large promise floor $\lambda\gg\epsilon$ can
favor the floor factory.

\begin{corollary}[Centered normalization]
\label{cor:gelulayernorm}
For $\epsilon>0$, $\mu=n^{-1}\sum_j a_j$, and $|a_j|\le R$, regularized centered
normalization
\[
 \frac{a_i-\mu}{\sqrt{\epsilon+n^{-1}\sum_j(a_j-\mu)^2}}
\]
has an exact sign factory at output radius $4\sqrt n$, using at most
$1+4nR^2/\epsilon$ original source signs in expectation.
\end{corollary}
\paragraph{Proof idea.}
Centering is a signed mixture of one coordinate and the uniform
coordinate average. It enlarges the coordinate radius by a factor of two,
after which the Gaussian normalization theorem applies directly.

\begin{proof}
A sign of mean $(a_i-\mu)/(2R)$ is a fair choice between an $a_i/R$
sign and the negative of a sign from a uniformly chosen coordinate.
It costs one original request. Apply Theorem~\ref{thm:gaussianrms}
to the centered vector with coordinate bound $2R$.
\end{proof}

Thinning trades output radius for source cost. The next corollary
opens a known gate before any Gaussian work and reaches radius
$2R/\sqrt\epsilon$, which is the radius of the floor factory at
$\lambda=\epsilon$.

\begin{corollary}[Thinning the vector Gaussian interface]
\label{cor:gaussian-large-radius}
Under Theorem~\ref{thm:gaussianrms}, put $\kappa=R^2/\epsilon$ and
suppose $\kappa\ge4n$.  There is an exact factory returning the
normalized coordinate at radius $2\sqrt\kappa$ with at most
\[
 2\sqrt{n/\kappa}\,(1+n\kappa)
\]
expected source requests and
\[
 O\!\left(B^4[1+n^{3/2}\sqrt\kappa]\right)
\]
expected auxiliary bit work.
\end{corollary}
\paragraph{Proof idea.}
Open the radius-adjustment gate before drawing any Gaussians.
Both the source count and the expensive local work are then charged only
on the gate's open branch.

\begin{proof}
Before doing any Gaussian work, open a known gate with probability
$p=2\sqrt{n/\kappa}\le1$.  On its open branch run the Gaussian
factory at radius $4\sqrt n$, and otherwise return a fair sign.
The output radius becomes $(4\sqrt n)/p=2\sqrt\kappa$.
The source bound follows by multiplication.
Lemma~\ref{lem:gaussianrmsbits} gives conditional auxiliary cost
$O(n(1+\kappa)B^4)$.  Multiplying by $p$, and using
$\kappa\ge4n\ge4$, gives the stated bound after adding the
$O(B^4)$ initial gate.  As before, every gate is sampled by certified
interval comparisons.
\end{proof}

\subsubsection{Finite-bit implementation}
\label{sec:gaussian-rms-bits}

\begin{lemma}[Linear stabilizer dependence in bit work]
\label{lem:gaussianrmsbits}
Let $n\ge1$ and let $R,\epsilon>0$ be rational. Put
$\kappa=R^2/\epsilon$, and let $B\ge16$ include their binary encoding
lengths and $\lceil\log_2(n+2)\rceil$. The factory in
Theorem~\ref{thm:gaussianrms}, including all Gaussian draws, weighted
coordinate selections, Poisson variables, random bits, counters, and
precision refinement, has expected auxiliary bit work
\begin{equation}
 O\!\left(n(1+\kappa)B^4\right).
 \label{eq:gaussianrmsbits}
\end{equation}
Source-sign invocations are charged separately. The implicit constant
is absolute. The output radius $4\sqrt n$ may be replaced by the
integer radius $4\lceil\sqrt n\rceil$ without increasing the expected
source count or changing this bit bound.
\end{lemma}
\paragraph{Proof idea.}
Two decompositions control precision costs. Split the Poisson mean
into a rational bulk and a bounded residual, and select Gaussian coordinates
from an integer envelope. Only a selected coordinate needs further
refinement; Gaussian fourth moments then keep the stabilizer dependence
linear.

\begin{proof}
We describe rational interval algorithms and bound their refinement
costs. Elementary series and their arithmetic are evaluated in dyadic
fixed-point intervals with outward rounding; logarithmically many guard
bits in the number of summed terms bound the accumulated rounding error.
Thus denominators from separate series terms are never multiplied together.
Each comparison uses an independent uniform bit stream. Almost
sure decisions suffice; equality with a random continuously distributed
comparison variable has probability zero.

\paragraph{Gaussian intervals.}
Two independent uniform variables $U,V\in(0,1)$ give independent standard
Gaussians by
\[
 G_1=\sqrt{-2\log U}\cos(2\pi V),\qquad
 G_2=\sqrt{-2\log U}\sin(2\pi V).
\]
The change to polar coordinates proves this distribution. The first
nonzero bit of $U$ locates it in $[2^{-L},2^{1-L}]$. The integer $L$
has a geometric tail. The proof of Lemma~\ref{lem:productaritybits}
gives certified intervals for $\pi$ and, after binary range reduction,
for $\log U$, at a bit cost cubic in the precision. The
usual sine and cosine series have factorial tail bounds on their
bounded arguments. Square-root intervals can be found by rational
bisection or integer square roots. Using twice the desired absolute
precision before the square-root step handles arbitrarily small
nonnegative arguments.
These procedures produce an interval of width at most $2^{-q}$ for a
Gaussian in $O((q+L+1)^3)$ bit operations. They never evaluate a
Gaussian as an exact real.

Generate $G_1,\ldots,G_n$ and the additional Gaussian $Z$ used by the
erf factory in this way. Set $S=\sum_j|G_j|$. A dyadic interval locating
$S$ within a factor of two is obtained by refinement. Its precision has
an exponential tail after $O(\log(n+2))$ initial bits: if this search has
not resolved at absolute precision $2^{-q}$, then
$S\le Cn2^{-q}$ for an absolute constant $C$, whereas
\[
 \Pr(S\le u)\le\Pr(|G_1|\le u)\le u
 \quad(u\ge0).
\]
The last bound uses the standard normal density and may be weakened to
$2u$ without changing any estimate. The signs of the finitely many
Gaussians are found by the same refinement; the union bound gives
$\Pr(\min_j|G_j|\le u)\le2nu$.

To simplify the remaining accounting, let $H$ be one plus the largest
leading-zero count in the Gaussian uniforms, the largest required
logarithmic precision for these signs and the dyadic scale of $S$, and
the logarithmic magnitude of the Gaussian variables. It can be
increased by an absolute multiple of $\log(n+2)$. The preceding bounds,
the normal density bound, and the Gaussian tail give, for every fixed
integer $r\ge1$,
\begin{equation}
 \mathbb E H^r\le C_r[1+\log(n+2)]^r.
 \label{eq:gaussianHmoments}
\end{equation}
The constants $C_r$ are absolute. A union bound over $n+1$ variables
followed by integration of an exponential tail proves this display.

\paragraph{The erf gate.}
The gate compares $|Z|$ with $\sqrt2 c$, where
$c=RS/\sqrt{2n\epsilon}$. At refinement level $q$, evaluate both sides
in intervals of absolute width $2^{-q-3}$. The arithmetic cost is
$O(n(B+H+q)^3)$. Conditional on $G_1,\ldots,G_n$, the variable $Z$ is
independent and its density is bounded. Therefore the probability of
an unresolved gate at level $q$ is at most $C2^{-q}$, uniformly in the
threshold $c$. Include the gate's extra refinement depth in $H$.
This leaves \eqref{eq:gaussianHmoments} valid with different absolute
constants. The event of exact equality has probability zero.

\paragraph{A Poisson variable with a random computable mean.}
On the open gate branch, $\mu=c^2-Z^2/2\ge0$. First compute a rational
interval $[\ell_0,u_0]$ for $\mu$ of width at most $1/4$. Set
$\ell=\max\{0,\ell_0\}$ and $r=\mu-\ell$. Then
\[
 0\le\ell\le\mu\le c^2,\qquad 0\le r\le1/4.
\]
No decision about the fractional part of $\mu$ is needed. The lower
approximation $\ell$ is an explicitly represented rational. Generate
$P_\ell\sim\operatorname{Pois}(\ell)$ from $\lfloor\ell\rfloor$
independent Poisson variables of mean one and one with the known rational
mean $\ell-\lfloor\ell\rfloor$. This part uses no further Gaussian
refinement. Independently generate $P_r\sim\operatorname{Pois}(r)$ by
thinning a Poisson variable of mean one with success probability $r$.
Each thinning test compares a fresh uniform with successive intervals
for $r=\mu-\ell$. The retained count is exactly Poisson with mean $r$,
so $P_\ell+P_r$ has the required Poisson law.

Lemma~\ref{lem:smallpoissonattention} generates a Poisson variable
whose mean is a rational in $[0,1]$ of bit length $b$ with $O(b^4)$
expected bit work and finite polynomial work moments; at mean one this
work is constant. The bulk rational part therefore costs
$O((1+c^2)(B+H)^4)$ conditionally on the Gaussian variables. The number
of Gaussian-dependent thinning comparisons has mean one. Each comparison
has a geometric refinement tail, because its uniform is fresh, and
evaluating $r$ to $q$ bits costs $O(n(B+H+q)^3)$. This part costs
$O(n(B+H)^4)$ in conditional expectation. The denominator, scale, and
polynomial coefficients involved have logarithmic size $O(B+H)$;
in particular $B$ includes $\log(2+\kappa)$.

Unbounded Poisson counters do not change these estimates. For every
fixed $p$, a Poisson variable $P$ of mean $\lambda\ge1$ satisfies
\[
 \mathbb E[(1+P)(1+\log(P+2))^p]
 \le C_p(1+\lambda)(1+\log(\lambda+2))^p.
\]
One proof divides the values into $P\le4\lambda$ and the intervals
$2^j\lambda<P\le2^{j+1}\lambda$, and uses
$\Pr(P\ge v)\le(e\lambda/v)^v$ for $v\ge\lambda$. Summing the
resulting convergent series proves the display. It accounts for
counters, majority arities, and source-address bookkeeping. The case
$\lambda\le1$ has an absolute bound by the same argument.

\paragraph{Choosing a Gaussian-weighted coordinate.}
The Gaussian vector is fixed throughout one erf call. Find a positive
rational $\sigma$ for which $\sigma\le S\le2\sigma$, using the interval
scale search described above. Choose a dyadic number $h$ with
\[
 \frac{\sigma}{32n}\le h\le\frac{\sigma}{16n}.
\]
For each $j$, compute an interval for $|G_j|$ of width at most $h/4$,
and let $b_j$ be its rational upper endpoint. Define
\[
 u_j=h\max\{1,\lceil b_j/h\rceil\},\qquad U=\sum_j u_j.
\]
Only rational rounding is used. The inequalities
\[
 |G_j|\le u_j\le |G_j|+2h,\qquad u_j\ge h,
 \qquad S\le U\le S+\sigma/8\le9S/8
\]
follow directly. All weights $u_j/h$ are positive integers, and their
sum is at most $66n$, hence also at most $72n$. Prepare their integer
prefix table once. This costs $O(n(B+H)^3)$ bits of work.

For every required coordinate draw, propose $j$ with probability $u_j/U$
using the integer table and retain it with probability $|G_j|/u_j$.
At acceptance, its law is exactly proportional to $|G_j|$. The success
probability is $S/U\ge8/9$, so the mean number of proposals is at most
$9/8$. A rejection comparison refines only the selected Gaussian.
Since $u_j\ge h$, obtaining absolute error $2^{-q}$ in this ratio costs
$O((B+H+q)^3)$, without a factor $n$. Fresh comparison uniforms give a
geometric refinement tail. The conditional expected work per accepted
coordinate is therefore $O((B+H)^4)$, including its address and sign.
The majority arity is drawn before its input signs are requested, and
has conditional mean at most $1+2c^2$. Predictable charging bounds the
expected total coordinate-selection work by
$O((1+c^2)(B+H)^4)$ in addition to the one-time table preparation.

\paragraph{Averaging the random precision cost.}
The preceding bounds, including the initial Gaussian interval work,
are controlled by
$C\,\mathbb E[(n+1+c^2)(B+H)^4]$. This expectation remains linear in
$\kappa$. In fact, $S^4\le n^3\sum_jG_j^4$, so
\[
 \mathbb E c^4
 =\frac{\kappa^2}{4n^2}\mathbb E S^4
 \le\frac34n^2\kappa^2,
 \qquad
 2\mathbb E c^2\le n\kappa.
\]
Consequently
$\mathbb E(1+c^2)^2\le(1+n\kappa)^2$. By the triangle inequality in
$L^2$, $\{\mathbb E(n+1+c^2)^2\}^{1/2}\le n+1+n\kappa$.
Cauchy--Schwarz and \eqref{eq:gaussianHmoments}, with $r=8$, give
\[
 \mathbb E[(n+1+c^2)(B+H)^4]
 \le(n+1+n\kappa)
       \{\mathbb E(B+H)^8\}^{1/2}
 \le Cn(1+\kappa)B^4.
\]
This also accounts for correlations between large Gaussian values,
large Poisson means, and difficult comparisons. The constants do not
depend on $R,\epsilon,n$, or the unknown source means.

Finally, the sine factory draws its odd degree before any Gaussian
factory calls. Its degree has an exponential tail and mean less than
one; all its scalar parameters are fixed computable constants. A
known constant gate and a bounded-parameter Poisson generator give
$O(B^4)$ auxiliary work. Conditional charging proves
\eqref{eq:gaussianrmsbits}. Every comparison terminates almost surely,
every Poisson variable is finite almost surely, and every selected
majority is finite. The entire factory therefore terminates almost
surely.

Let $m=\lceil\sqrt n\rceil$, found by integer arithmetic. Thin the
output by the known algebraic probability $\sqrt n/m$, returning a
fair sign otherwise. Squaring rational interval endpoints gives an
exact implementation of this comparison with $O(B^4)$ expected work.
The new mean is the required normalized coordinate at radius $4m$.
\end{proof}
